\originalchapter{顶点染色}\label{ch:coloring}
\sourceof{18.315 L4--5、L8; 删双顶点引理及换色过程为编者补证}
\originalsection{染色数与贪心排序}
\begin{definition}
正常顶点染色是映射 $c:V\to\{1,\ldots,k\}$, 使每条边 $uv$ 满足 $c(u)\ne c(v)$. 最小所需颜色数为染色数 $\chi(G)$. 团是内部两两相邻的顶点集, 最大大小记 $\omega(G)$; 最大独立集大小记 $\alpha(G)$.
\end{definition}
染色把顶点分成若干独立集. 团内各点必须异色, 每个颜色类至多含 $\alpha$ 点, 所以 $\chi\ge\omega$, 且对 $n>0$ 有 $\chi\ge n/\alpha$. 这些界不总精确: 奇圈团数为2, 染色数却为3.
\begin{proposition}\label{prop:greedycolor}
有限简单图满足 $\chi(G)\le\Delta+1$. 更一般地, 若每个非空子图有度数至多 $k$ 的顶点, 则 $\chi(G)\le k+1$.
\end{proposition}
\begin{proof}
依任意顺序染点, 一个顶点的已染邻点至多 $\Delta$ 个, 总有一种未被使用的颜色. 对第二结论, 反复在剩余诱导子图删除一个度至多 $k$ 的顶点, 再按删除顺序的逆序染色. 每点染色时至多有 $k$ 个已染邻点.
\end{proof}
第二种图称为 $k$ 退化图. 它控制的是每个子图的稀疏性, 而不是全图平均度. 一个大图的平均度很小, 仍可能含一个巨大团; 只检查全图不能直接给出退化染色界.
\begin{proposition}
若 $m$ 条边的图需要 $q$ 种颜色, 则 $m\ge\binom q2$, 因而 $\chi(G)\le (1+\sqrt{1+8m})/2$.
\end{proposition}
\begin{proof}
在一个最少颜色的染色中, 任意两个颜色类之间必须有边; 否则可把两类合并而少用一色. 不同颜色类对需要不同的边, 所以至少有 $\binom q2$ 条边. 解二次不等式即得上界.
\end{proof}
\originalsection{Brooks 定理}
\begin{lemma}\label{lem:lowdegreecolor}
连通图 $F$ 的最大度不超过 $D$. 若有顶点 $r$ 的度至多 $D-1$, 则 $F$ 可用 $D$ 色正常染色.
\end{lemma}
\begin{proof}
取以 $r$ 为根的生成树, 按深度从大到小染色, 最后染根. 每个非根顶点至少有一个未染的父顶点, 所以已染邻点至多 $D-1$ 个. 根的总度也至多 $D-1$, 贪心过程能完成.
\end{proof}
\begin{lemma}\label{lem:brooksorder}
若 $G$ 为 $D$ 正则、二连通、非完全的简单图, 且 $D\ge3$, 则有非相邻顶点 $a,b$ 及其公共邻点 $v$, 使 $G-\{a,b\}$ 连通.
\end{lemma}
\begin{proof}
若 $G$ 三连通, 选一条非相邻顶点间的最短路, 它的前三点构成诱导路 $a,v,b$. 三连通性直接保证删两端后仍连通.

否则有二点割 $\{v,c\}$. 图 $H=G-v$ 连通, 而 $c$ 是其割点. 由块树引理\ref{lem:blocktree}, $H$ 有两个不同叶块. 每个叶块的内部点集中必须有 $v$ 的邻点: 若没有, 删去该块唯一的连接割点就会把内部点与 $v$ 分离, 违反 $G$ 二连通. 在这两个内部点集中分别取 $a,b\in N_G(v)$. 它们处于不同块的内部, 不相邻. 删去 $a,b$ 后, 每个叶块的剩余部分仍连到原连接点, 故 $H-a-b$ 连通. 由于 $v$ 原有 $D\ge3$ 个邻点, 删除 $a,b$ 后仍有邻点, 所以 $G-a-b$ 也连通.
\end{proof}
\begin{theorem}[Brooks]\label{thm:brooks}
有限简单连通图若既不是完全图, 也不是奇圈, 则 $\chi(G)\le\Delta(G)$.
\end{theorem}
\begin{proof}
最大度不超过2时, 连通图为路径、圈或单点. 排除完全图和奇圈后是非完全路径或偶圈, 都可用2色染色. 以下设 $D=\Delta\ge3$. 如果图非正则, 用低度引理. 若有割点 $v$, 对 $G-v$ 的每个分支 $C_i$ 考虑连通图 $G[C_i\cup\{v\}]$. 它最大度不超过 $D$, 且 $v$ 在其中的度至多 $D-1$, 因为其他分支至少还占有一个邻点. 各部分都可用 $D$ 色染色, 置换各部分颜色使 $v$ 同色后即可拼接.

剩下正则二连通情形, 用上一引理找 $a,v,b$. 先把不相邻的 $a,b$ 同染颜色1. 在 $G-a-b$ 中取以 $v$ 为根的生成树, 从最深点向根贪心染色. 每个非根点有未染父顶点, 可用 $D$ 色. 最后根 $v$ 的两个邻点 $a,b$ 同色, 全部邻点使用的不同颜色至多 $D-1$, 所以根也有颜色可用.
\end{proof}
定理指出 $\Delta+1$ 的例外恰由完全图和奇圈产生. 它没有说任意给定染色都能不加颜色局部改成最优染色; 证明中构造的顶点顺序是实质部分.
\originalsection{Mycielski 构造}
团数与染色数之间的差距可以任意大, 即使禁止三角形也不能仅凭团数控制染色数.
\begin{theorem}
若无三角形图 $G$ 有 $\chi(G)=k\ge2$, 则有无三角形图 $M(G)$ 满足 $\chi(M(G))=k+1$.
\end{theorem}
\begin{proof}
设原顶点为 $v_1,\ldots,v_n$, 加副本 $u_1,\ldots,u_n$ 及一点 $w$. 保留原边; 对每条原边 $v_iv_j$, 加 $u_iv_j,u_jv_i$; 最后令 $w$ 邻接所有 $u_i$. 副本之间无边, $w$ 不邻接原点. 若三角形含 $w$, 另外两点须为两个副本而不相邻, 不可能. 若含一个副本 $u_i$, 另外两原点与 $v_i$ 形成原三角形, 也不可能; 其余三角形都在原图中. 所以新图无三角形.

沿用原 $k$ 色染色, 给 $u_i$ 与 $v_i$ 同色, 给 $w$ 新色, 得到 $k+1$ 色. 假如只需 $k$ 色, 令 $w$ 的颜色为 $k$, 则全部 $u_i$ 都不取 $k$. 在原图中, 将每个颜色 $k$ 的 $v_i$ 改为 $u_i$ 的颜色. 两个被改点原先同色, 本来不相邻; 被改点 $v_i$ 与未改邻点 $v_j$ 的新色不同, 因为新图有边 $u_iv_j$. 得原图的 $k-1$ 色染色, 矛盾.
\end{proof}
从 $C_5$ 开始反复构造, 可得任意高染色数的无三角形图. 构造和第\ref{ch:ramsey}章的随机方法分别展示确定性与概率性两种存在证明.
\originalsection{染色的重新配置}
\begin{theorem}\label{thm:recolor}
图 $G$ 最大度为 $D$, 若 $k\ge D+2$, 则任意两个正常 $k$ 色染色之间, 可通过每次只改变一个顶点的颜色、全程保持正常染色而互相转换.
\end{theorem}
\begin{proof}
对顶点数归纳. 删除一点 $v$, 归纳把余图从初始染色变为目标染色. 模拟余图每次把一点 $u$ 改为颜色 $b$ 的操作. 若 $u$ 不邻接 $v$ 或 $v$ 当前颜色不是 $b$, 直接执行. 否则先给 $v$ 换成一种颜色, 既不同于所有当前邻点颜色, 又不同于 $b$. 被禁颜色至多 $D+1$ 种, 在 $k\ge D+2$ 时总能找到; 随后执行 $u$ 的变化. 余图全部达到目标后, 把 $v$ 改成目标色, 这时所有邻点已是目标色, 不发生冲突.
\end{proof}
颜色数充足保证状态图连通, 不自动给出快速随机采样或最短换色路径. 18.315 L8 还讨论格点三染色的高度函数及直径; 那个专门模型不在这里被泛化为所有图的同一结论.
\originalsection{习题与解答}
\begin{exercise}\label{ex:col1}
证明任何森林可用两色染色, 并给出一个最大度任意大但染色数为2的连通图.
\end{exercise}
\begin{exercise}\label{ex:col2}
证明若 $G$ 的染色数为 $k\ge1$, 则它含有最小度至少 $k-1$ 的子图.
\end{exercise}
习题\ref{ex:col1}的解答.
\begin{solution}
在每棵树中按根的距离奇偶染色; 森林无奇圈, 也可直接用二部判据. 星 $K_{1,D}$ 在 $D\ge1$ 时最大度 $D$ 而染色数2, 说明 Brooks 给出上界而不是精确公式.
\end{solution}
习题\ref{ex:col2}的解答.
\begin{solution}
选顶点数最少的诱导子图 $H$, 使 $\chi(H)=k$. 若某点度至多 $k-2$, 删它后由最小性可用至多 $k-1$ 色染色, 再给它选未被邻点占用的一色, 得 $H$ 的 $k-1$ 色染色, 矛盾. 所以 $\delta(H)\ge k-1$.
\end{solution}
